% ICSE 2027 Research Track evaluation chapter.
% Embeddable IEEE conference fragment; not a standalone document.
% Required packages: graphicx, amsmath, booktabs, tabularx, placeins,
% and dblfloatfix.

% Overleaf resolves image paths relative to the selected main document, not
% relative to this included fragment.  This helper accepts the repository
% layout, a conventional Overleaf figures/ folder, same-folder uploads, and
% PNG fallbacks.
\providecommand{\evalfigure}[2][]{%
  \IfFileExists{figures/#2.pdf}{\includegraphics[#1]{figures/#2.pdf}}{%
  \IfFileExists{docs/figures/#2.pdf}{\includegraphics[#1]{docs/figures/#2.pdf}}{%
  \IfFileExists{#2.pdf}{\includegraphics[#1]{#2.pdf}}{%
  \IfFileExists{figures/#2.png}{\includegraphics[#1]{figures/#2.png}}{%
  \IfFileExists{docs/figures/#2.png}{\includegraphics[#1]{docs/figures/#2.png}}{%
  \fbox{\parbox{0.9\columnwidth}{\centering Missing figure: #2}}}}}}}}

\section{Evaluation}
\label{sec:evaluation}

We evaluate \emph{EsCapturer-full} through four research questions that map
directly to its empirical claims:
\begin{itemize}
\setlength{\itemsep}{2pt}
\setlength{\parsep}{0pt}
\setlength{\topsep}{3pt}
\item \textbf{RQ1: Detection Effectiveness and Split Reliability.} How
accurately does EsCapturer-full detect malicious traces, and how does its
performance change under exact- and source-group-disjoint partitioning?
\item \textbf{RQ2: Component Attribution and Routing Behavior.} What do the
semantic elements, graph view, and interleaving prior contribute, and do the
learned routing weights follow the intended structural prior?
\item \textbf{RQ3: Robustness.} How stable is a clean-trained detector under
API/syscall insertion, deletion, and bounded local reordering?
\item \textbf{RQ4: Efficiency and Deployment Practicality.} What computation,
memory, and deployment-time overhead does the complete pipeline introduce?
\end{itemize}

\subsection{Experimental Setup}
\label{sec:eval-setup}

\subsubsection{Datasets and Split Protocol}
Quo Vadis contains report-level ordered Windows API sequences generated by
Speakeasy, with each report capped at 5,000 calls.  Zenodo record 11079764
contains VMI API traces divided into windows of at most 200 calls with stride
100.  For the balanced benchmark, each dataset contributes 50,000 records
(25,000 benign and 25,000 malware) and uses a seed-7
35,000/7,500/7,500 train/validation/test split.

Random record-level splitting can place duplicate sequences or related source
groups in different partitions.  We therefore construct a stricter audit by
removing exact sequences with conflicting labels, collapsing remaining exact
duplicates, and assigning report/SHA groups atomically to a 70/15/15 split.
Table~\ref{tab:dataset-protocol} summarizes the resulting partitions; both
exact-sequence and source-group overlap are zero.

\begin{table}[!b]
\caption{Balanced benchmark and leakage-safe audit construction.  Removed
counts refer to records.}
\label{tab:dataset-protocol}
\centering
\scriptsize
\setlength{\tabcolsep}{3.5pt}
\begin{tabular}{lrrrr}
\hline
Dataset & Balanced & Conflict & Duplicate & Audit final \\
\hline
Quo Vadis & 50,000 & 6,895 & 36,469 & 6,636 \\
Zenodo 11079764 & 50,000 & 3,601 & 7,517 & 38,882 \\
\hline
\end{tabular}
\end{table}

\subsubsection{Baselines and Metrics}
All scores are produced by local runs under the same split and metric
definitions rather than copied from prior publications.  The comparison
covers frequency and $n$-gram classifiers, sequence models, graph models, and
adapted malware-detection systems: Nebula, API2Vec++, DawnGNN,
GPT4-API-BERT-CNN, DeepCapa, APILI, and MME-TextCNN.  Adapted and
reimplemented methods are identified explicitly because their original tasks
or released pipelines differ from binary runtime-trace detection.

We report Accuracy, Precision, Recall, F1, and ROC-AUC for the malicious
class.  Validation data alone determine checkpoints and decision thresholds;
the selected threshold is then fixed for test inference.  Three-seed audit
results report the arithmetic mean and standard deviation over seeds 7, 17,
and 27.  Single-seed balanced, ablation, robustness, and routing analyses are
presented as bounded evidence and do not support statistical-superiority
claims.

\subsubsection{Implementation}
EsCapturer-full uses frozen semantic templates, sequence and graph encoders,
and interleaving-prior adaptive fusion.  Vocabularies, caches, normalizers,
and graph statistics are fitted on training records only.  The matched Full
model and ablations use 32-dimensional embeddings, BCE loss, learning rate
$5{\times}10^{-4}$, at most 30 epochs, patience 5, maximum sequence length
512, and at most 16 uniformly covered behavior units.  The semantic extractor
is frozen and cached, so neither training nor deployment-time inference makes
online LLM calls.

\subsection{RQ1: Detection Effectiveness and Split Reliability}
\label{sec:rq1}

RQ1 evaluates the complete detector on the balanced benchmarks and the
stricter leakage-safe audit.  The former measures basic effectiveness under a
common binary split; the latter tests whether discrimination remains after
duplicate removal and source-group isolation.

\subsubsection{Balanced-Benchmark Performance}
Table~\ref{tab:rq1-balanced-main} reports F1 and AUC for the two independent
50,000-record benchmarks.  Figure~\ref{fig:rq1-balanced-paired} exposes the
gap between thresholded detection and threshold-independent ranking.

\begin{table*}[!t]
\caption{Local detection results on the two 50,000-record balanced
benchmarks.  All methods use the common binary task and seed-7 split.}
\label{tab:rq1-balanced-main}
\centering
\scriptsize
\setlength{\tabcolsep}{5pt}
\begin{tabular}{lrrrr}
\hline
Method & Quo F1 & Quo AUC & Zenodo F1 & Zenodo AUC \\
\hline
EsCapturer-full & 0.9626 & 0.9935 & 0.9409 & 0.9791 \\
TF-IDF/SGD prototype & 0.9618 & 0.9947 & 0.9516 & 0.9874 \\
Nebula & 0.9135 & 0.9708 & 0.7921 & 0.8994 \\
API2Vec++ & 0.9281 & 0.9824 & 0.9069 & 0.9611 \\
DawnGNN-reimpl & 0.8800 & 0.9458 & 0.8030 & 0.9189 \\
GPT4-API-BERT-CNN & 0.9544 & 0.9926 & 0.9228 & 0.9737 \\
DeepCapa-adapted & 0.9648 & 0.9953 & 0.9452 & 0.9863 \\
APILI-adapted & 0.9388 & 0.9869 & 0.9368 & 0.9806 \\
MME-TextCNN & 0.9389 & 0.9890 & 0.9287 & 0.9789 \\
\hline
\end{tabular}
\end{table*}

\begin{figure*}[!t]
\centering
\evalfigure[width=0.94\textwidth]{rq1_balanced_paired_scores}
\caption{Detection effectiveness on the balanced benchmarks.  Circles and
diamonds report F1 and AUC, respectively; EsCapturer-full is highlighted.}
\label{fig:rq1-balanced-paired}
\end{figure*}

On Quo Vadis, EsCapturer-full obtains Accuracy 0.9631, Precision 0.9743,
Recall 0.9512, F1 0.9626, and AUC 0.9935.  It has the highest precision among
the compared methods, while DeepCapa-adapted is 0.0022 and 0.0018 higher in
F1 and AUC.  On Zenodo, EsCapturer-full obtains Accuracy 0.9405, Precision
0.9345, Recall 0.9475, F1 0.9409, and AUC 0.9791.  TF-IDF/SGD and
DeepCapa-adapted are higher on several metrics.  The balanced benchmark thus
supports competitive effectiveness, not universal superiority.

\subsubsection{Leakage-Safe Audit}
The Quo Vadis audit contains 4,647/994/995 train/validation/test records.
Across three training seeds, EsCapturer-full obtains Accuracy
$0.8911{\pm}0.0065$, Precision $0.8453{\pm}0.0047$, Recall
$0.7021{\pm}0.0344$, F1 $0.7667{\pm}0.0191$, and AUC
$0.9168{\pm}0.0121$.  The complete audit comparison is reported in the
supplementary material rather than treated as a headline leaderboard.

A post-hoc overlap check finds that 90.0\% of records in the balanced Quo
Vadis test fold have an exact sequence match in its training fold.  The audit
removes this shortcut, reduces the training set from 35,000 to 4,647 records,
and changes the malware-family distribution.  Its F1 decrease is driven
primarily by Recall ($0.7021$) rather than Precision ($0.8453$), with the
largest miss rates on the smaller RAT, Coinminer, and Trojan cohorts.  Thus,
the audit measures generalization to unique source groups rather than merely
re-scoring the balanced population.

\noindent\textbf{Answer to RQ1.}
EsCapturer-full is competitive on both balanced datasets, but its performance
drops under exact- and source-group-disjoint evaluation.  The current evidence
supports effective record-level detection, while generalization to unique
source groups remains a limitation rather than a performance claim.

\subsection{RQ2: Component Attribution and Routing Behavior}
\label{sec:rq2}

RQ2 tests the three design elements most closely tied to the method:
semantic abstraction, structural modeling, and interleaving-conditioned
fusion.  Every matched variant uses the same balanced seed-7 split, training
budget, checkpoint rule, and validation-selected threshold as Full.

\subsubsection{Matched Ablation}
Figure~\ref{fig:rq2-ablation-delta} reports the absolute F1 change from Full.
On Quo Vadis, removing the graph view and interleaving prior decreases F1 by
0.0026 and 0.0030, whereas removing semantic elements increases F1 by 0.0007.
On Zenodo, removing semantic elements and the graph view decreases F1 by
0.0120 and 0.0128; removing the prior increases F1 by 0.0013.  AUC changes
follow the same broad pattern but remain below 0.0035 in magnitude.  The graph
view is the only component with a consistent positive direction across both
datasets.

\begin{figure}[!t]
\centering
\evalfigure[width=\columnwidth]{rq3_ablation_delta_f1}
\caption{Matched ablation effects relative to Full.  Bars report F1 change in
percentage points; positive values mean that removing the component improved
F1.}
\label{fig:rq2-ablation-delta}
\end{figure}

\subsubsection{Routing-Weight Behavior}
Table~\ref{tab:rq2-routing} aggregates graph-view gate weights over the two
balanced test folds.  The gate is effectively binary: non-interleaved pairs
receive near-zero graph weight and interleaved pairs receive near-one graph
weight.  This direction agrees with the encoded prior, but the saturation
provides little instance-specific nuance and is not causal evidence that the
prior improves detection.  The Zenodo ablation, where removing the prior
slightly improves F1, reinforces this distinction.

\begin{table}[!t]
\caption{Graph-view routing weights on the balanced test folds.}
\label{tab:rq2-routing}
\centering
\scriptsize
\setlength{\tabcolsep}{3pt}
\begin{tabular}{llrrr}
\hline
Dataset & Pair group & Pairs & Mean & Median \\
\hline
Quo Vadis & $C_{ij}=0$ & 878,022 & $<10^{-9}$ & $<10^{-9}$ \\
Quo Vadis & $C_{ij}=1$ & 747,408 & $>0.9999$ & 1.0000 \\
Zenodo & $C_{ij}=0$ & 383,856 & $<10^{-9}$ & $<10^{-9}$ \\
Zenodo & $C_{ij}=1$ & 1,536,144 & $>0.9999$ & 1.0000 \\
\hline
\end{tabular}
\end{table}

\noindent\textbf{Answer to RQ2.}
The graph view contributes consistently, while semantic abstraction and the
interleaving prior have dataset-dependent marginal effects.  Routing weights
follow the intended interleaving division but are nearly saturated, limiting
their value as fine-grained explanations.

\subsection{RQ3: Robustness}
\label{sec:rq3}

RQ3 perturbs the balanced test folds while keeping the seed-7 training data,
checkpoint, and threshold unchanged.  Insertion adds calls sampled from the
training-set benign-call distribution; deletion removes calls uniformly while
retaining at least one call; local reordering permutes calls within selected
five-call windows.  Each deterministic perturbation is applied at 10\%,
20\%, and 30\% intensity.  Figure~\ref{fig:rq3-robustness} reports the
resulting F1 loss.

\begin{figure}[!t]
\centering
\evalfigure[width=\columnwidth]{rq4_robustness_heatmap}
\caption{F1 loss under clean-trained robustness tests (percentage points).
Darker cells indicate larger degradation; every condition reuses the clean
checkpoint and validation-selected threshold.}
\label{fig:rq3-robustness}
\end{figure}

Insertion is the most damaging perturbation.  At 30\%, F1 decreases by
0.3096 on Quo Vadis and 0.2181 on Zenodo, with AUC falling to 0.8429 and
0.7635.  Deletion has an intermediate effect: at 30\%, F1 decreases by
0.1407/0.0923 and AUC remains 0.9397/0.9071 on Quo Vadis/Zenodo.  Local
reordering is least harmful; at 30\%, F1 remains 0.9433/0.9171 and AUC
remains 0.9890/0.9636.

The endpoint precision/recall values clarify the failure mode.  On Quo Vadis,
30\% insertion retains Precision 0.9309 but lowers Recall to 0.5029,
indicating missed detections.  On Zenodo, Precision/Recall fall to
0.7074/0.7389, indicating a broader loss of separation.

\noindent\textbf{Answer to RQ3.}
EsCapturer-full is stable to bounded local reordering and moderately robust to
deletion, but benign-call insertion increasingly degrades both thresholded
detection and ranking.

\subsection{RQ4: Efficiency and Deployment Practicality}
\label{sec:rq4}

RQ4 uses timing and memory counters stored by the two balanced Full runs.
Both profiles cover CPU inference over all 7,500 test records.  Semantic
extraction and graph construction are local preprocessing stages, and the
frozen template cache eliminates online LLM calls.

\begin{table}[!t]
\caption{Single-machine efficiency measurements.}
\label{tab:rq4-efficiency}
\centering
\scriptsize
\setlength{\tabcolsep}{4.5pt}
\begin{tabular}{lrr}
\hline
Measurement & Quo Vadis & Zenodo \\
\hline
Training time & 9.35 h & 10.03 h \\
Semantic extraction & 42.02 s & 28.91 s \\
Behavior-unit construction & 14.35 s & 11.22 s \\
Graph construction & 64.94 s & 56.04 s \\
Test latency & 13.41 ms/record & 15.04 ms/record \\
Test throughput & 74.6 records/s & 66.5 records/s \\
Parameters & 300,405 & 128,981 \\
Peak RSS & 1,718.7 MiB & 1,687.0 MiB \\
Checkpoint size & 1.30 MiB & 0.54 MiB \\
\hline
\end{tabular}
\end{table}

Inference remains below 16 ms per record, corresponding to 66--75 records/s.
Training requires roughly 9--10 CPU hours but is an offline operation.
Deployment can reuse the frozen 600.7 KiB Quo Vadis or 13.9 KiB Zenodo
template cache.  These measurements establish scale rather than
hardware-independent latency bounds.

\noindent\textbf{Answer to RQ4.}
The measured CPU implementation supports tens of records per second without
online LLM calls.  Its principal cost is offline training; model size and
per-record inference remain modest in the recorded environment.

\subsection{Threats to Validity}
\label{sec:eval-threats}

\textbf{Construct validity.}
The balanced samples do not preserve source prevalence, and the compact traces
contain API names without the arguments, return values, or runtime objects
needed to exercise every semantic field.  The synthetic perturbations model
padding, omission, and bounded order noise, but not adaptive attacks, argument
mutation, or object obfuscation.

\textbf{Internal validity.}
Duplicate sequences, related windows, and shared source identities can inflate
record-level results.  The leakage-safe audit removes conflicting and duplicate
sequences and enforces source-group-disjoint partitions.  Adapted baselines may
still differ from their original systems, so we identify them explicitly and
treat all reported values as local results.

\textbf{Conclusion validity.}
Balanced, ablation, robustness, and routing results use one seed-7 split.
Only the Quo Vadis leakage-safe audit uses three training seeds.  We therefore
avoid significance and universal-superiority claims; small ablation changes
should be interpreted as directional rather than definitive.

\textbf{External validity.}
Both datasets contain Windows API traces from specific analysis environments.
Results may not transfer to other sandbox versions, live endpoints, Linux
system calls, packed malware, future families, or different hardware.

\subsection{Summary of Findings}
\label{sec:eval-summary}

EsCapturer-full is competitive on the record-level benchmarks; the stricter
audit identifies generalization to unique source groups as an open limitation.
The graph view contributes consistently, whereas semantic abstraction and the
interleaving prior have dataset-dependent effects; the routing gate follows
the prior but saturates.  The detector is robust to local reordering, less
robust to deletion, and most sensitive to benign-call insertion.  Finally,
the recorded CPU implementation processes 66--75 records/s without online
LLM calls.

\FloatBarrier
